% ChassisModule
\section{【応用】ChassisModule — 機能統合パターン}

\begin{patternbox}{ChassisModule で学べるパターン}
\begin{itemize}
    \item 複数モジュールを内部に保持する「機能統合系モジュール」
    \item 内部モジュールのinit()/updateOutput()/deinit()を親が管理
    \item ベクトル合成による高レベル制御（3軸入力→4輪出力）
    \item 内部モジュールのSystemData登録が不要なケース
\end{itemize}
\end{patternbox}

\subsection{内部モジュールの保持}

\begin{lstlisting}[style=cppstyle, caption={ChassisModule.h — 内部にDriveMotorModuleを4つ保持}]
struct ChassisConfig {
    DriveMotorConfig motors[4];     // 4輪分のモーターConfig
    WheelPattern     wheelPattern;
    float            maxOutputPercent;
};

class ChassisModule : public IModule {
private:
    ChassisConfig      _config;
    DriveMotorModule   _motors[4];  // 内部にモジュールを直接保持
public:
    ChassisModule(const ChassisConfig& config);
    bool init() override;
    void updateOutput(SystemData& data) override;
    void deinit() override;
};
\end{lstlisting}

\begin{pointbox}[内部モジュールはIModule配列に登録しない]
\texttt{\_motors[4]}は\texttt{ChassisModule}のメンバ変数であり、
\texttt{main.cpp}の\texttt{outputModules[]}には登録しない。
\texttt{SystemData}にも\texttt{DriveMotorData}を含める必要がない。
外部からは\texttt{ChassisModule}1つだけが見え、3軸入力だけで4輪を制御できる。
\end{pointbox}

\subsection{親モジュールからの制御委譲}

\begin{lstlisting}[style=cppstyle, caption={ChassisModule.cpp — init()/deinit()で内部モジュールを管理}]
ChassisModule::ChassisModule(const ChassisConfig& config)
    : _config(config),
      _motors{
          DriveMotorModule(config.motors[0]),
          DriveMotorModule(config.motors[1]),
          DriveMotorModule(config.motors[2]),
          DriveMotorModule(config.motors[3]),
      } {}
\end{lstlisting}

\begin{notebox}[コンストラクタの初期化リスト — 初学者向け解説]
コロン（\texttt{:}）の後に書かれている部分が「初期化リスト」です。ここでは以下の処理を行っています。

\begin{enumerate}
    \item \texttt{\_config(config)} — Config構造体をコピー保存
    \item \texttt{\_motors\{...\}} — DriveMotorModuleの配列を、各モーターのConfigで初期化
\end{enumerate}

配列の初期化は波括弧 \texttt{\{\}} で囲みます。中の \texttt{DriveMotorModule(config.motors[0])} は「config.motors[0]の設定でDriveMotorModuleを作る」という意味です。

\vspace{2mm}
これは「コンストラクタの中で4つの子モジュールを一気に作る」処理です。
初期化リスト構文に慣れないうちは「コンストラクタ引数を受け取って、メンバ変数に渡している」とだけ理解すれば十分です。
\end{notebox}

\begin{lstlisting}[style=cppstyle, caption={ChassisModule.cpp — init()/deinit()（続き）}]

bool ChassisModule::init() {
    for (int i = 0; i < 4; i++) {
        if (!_motors[i].init()) return false;
    }
    return true;
}

void ChassisModule::deinit() {
    for (int i = 0; i < 4; i++) {
        _motors[i].deinit();
    }
}
\end{lstlisting}

\begin{lstlisting}[style=cppstyle, caption={ChassisModule.cpp — updateOutput()でベクトル合成→各モーターをdrive()}]
void ChassisModule::updateOutput(SystemData& data) {
    float outputs[4];
    // ベクトル合成（ホイールパターンに応じた方向行列）
    calcWheelOutputs(data.chassis, outputs);
    // リミッタ + スケーリング
    applyLimiter(outputs, _config.maxOutputPercent);
    // 各内部モジュールのdrive()を直接呼ぶ
    for (int i = 0; i < 4; i++) {
        _motors[i].drive(outputs[i]);
        data.chassis.motorOutputs[i] = outputs[i];
    }
}
\end{lstlisting}

\begin{pointbox}[機能統合の原則]
仕様書の「モジュール間の依存関係は存在しないものとして設計する」ルールに従い、
依存関係があるモジュール同士は1つに統合する。
ChassisModuleとDriveMotorModuleは明確に依存関係があるため、
ChassisModuleがDriveMotorModuleを内部部品として包含している。
\end{pointbox}

\begin{notebox}[Config構造体のネスト]
\texttt{ChassisConfig.motors[4]}は\texttt{DriveMotorConfig}の配列。
Config構造体は別のConfig構造体をメンバとして持てる。
これにより\texttt{ProjectConfig.h}で4輪分のピンアサインを一括定義できる。
\end{notebox}
